% Separately attributed quotient toolkit; original Hodges chapter closure is
% recorded outside the reader text in edition-2026-p03-imaginaries.md.
\section{Imaginaries and Quotient Interpretations}
\label{sec:p03-imaginaries-quotient-interpretations}

A definable quotient need not have a definable set of representatives.
An interpretation therefore specifies equality on representatives as well
as the domain and operations. This chapter gives that construction in arbitrary
finite dimensions and for many target sorts. The quotient account of
interpretation and translation is described by Hodges
\cite[\S2, paragraphs on interpretation]{Hodges2020ModelTheorySEP};
the quotient-sort and uniform-coding presentations are separately described in
\cite[Lecture 15, pp.~1--4]{Speirs2013Imaginaries} and
\cite[slides~6--8]{Baldwin2023Imaginaries}. The proofs below are explicit local
derivations of the declared presentation, not attributions of an unread
original proof to those sources.

\subsection{Definable domains, quotient equality, and translation}

\begin{hbconvention}{def:p03-imaginaries-standing}{Signatures, parameters, and dimensions}
The source language \(L\) and the interpreted language \(K\) are set-sized
first-order languages with finite-arity functions, relations, and logical
equality. Either may have many sorts; every home sort is nonempty.
Carriers retain their sort tags, and equality compares elements within one
sort only; identical underlying sets do not identify two sort symbols.
Let \(T\) be an \(L\)-theory. A finite tuple of parameters is handled by
naming its coordinates in \(L\), and by declaring the corresponding theory
in that expanded language before making a uniform claim.

For every \(K\)-sort \(s\), choose a finite, possibly empty, list of
\(L\)-sorts and a variable block \(\bar x_s\) of those sorts. Its length
\(n_s\) is the dimension, not the number of quotient classes. The empty
product contains its one empty tuple. Quantification over a block of length
zero performs no quantifier operation. All representative variables introduced
below are fresh, and bound variables are renamed before composition.
\end{hbconvention}

\begin{hbdefinition}{def:p03-imaginaries-interpretation}{A many-sorted definable quotient interpretation}
An interpretation \(I\) specifies, for each sort \(s\), formulas
\(D_s(\bar x_s)\) and \(E_s(\bar x_s,\bar y_s)\). Require
\(T\models\exists\bar x_s D_s\), and require \(T\) to imply that
\(E_s\) is an equivalence relation on \(D_s\).
For \(M\models T\), the interpreted sort is
\[
             I(M)_s=D_s(M)/E_s(M).
\]
Write \(q_s:D_s(M)\to I(M)_s\) for the set-theoretic quotient map.
It is not a chosen section or a new partial term of \(L\).

Each \(K\)-relation is assigned an \(L\)-formula on the corresponding
representative blocks, invariant under replacing any block by an equivalent
one. A nullary relation is assigned a sentence. For each function
\(f:s_1\times\cdots\times s_k\to t\), specify a graph formula
\(G_f(\bar x_1,\ldots,\bar x_k,\bar z)\), invariant in every block,
such that, on domain representatives, an output in \(D_t\) exists and any
two outputs are \(E_t\)-equivalent. Constants use the same condition with
\(k=0\). These conditions are required consequences of \(T\).
The graph is thus total and single-valued on quotient classes, not necessarily
single-valued on representatives. If \(U\) is the target theory, also require
\(T\models\sigma^I\) for each axiom \(\sigma\) of \(U\), using the
translation below.
\index{interpretation!definable quotient}\index{quotient equality}
\end{hbdefinition}

\begin{hbdefinition}{def:p03-imaginaries-translation}{Term values and formula pullback}
For a \(K\)-term \(t\) of sort \(s\), define
\(V_t^I(\bar x,\bar z)\), for \(\bar z\in D_s\), recursively.
A variable uses \(E_s\) between its input block and \(\bar z\);
a constant uses its graph. For a compound term use
\[
 V_{f(t_1,\ldots,t_k)}^I(\bar x,\bar z)
 =\exists\bar u_1\cdots\bar u_k
 \left(\bigwedge_{i=1}^k
       (D_{s_i}(\bar u_i)\land V_{t_i}^I(\bar x,\bar u_i))
       \land G_f(\bar u_1,\ldots,\bar u_k,\bar z)\right).
\]
An atomic relation \(R(t_1,\ldots,t_k)\) is translated by existentially
choosing domain representatives for each term value and applying the formula
assigned to \(R\). For equality \(t=u\), apply the assigned \(E_s\)
to the two chosen values. Translation commutes with Boolean connectives and
\[
 (\exists v_s\,\varphi)^I
   =\exists\bar x_s(D_s(\bar x_s)\land\varphi^I),\qquad
 (\forall v_s\,\varphi)^I
   =\forall\bar x_s(D_s(\bar x_s)\to\varphi^I).
\]
The variables free in \(\varphi^I\) are representative blocks for the
free target variables. Their domain conditions are hypotheses on assignments,
not silently appended axioms of the source structure.
\end{hbdefinition}

\begin{hbtheorem}{thm:p03-imaginaries-quotient-truth}{proved=annex:p03-imaginaries-quotient-truth}{Well-defined quotient structures and the truth lemma}
The data of \cref{def:p03-imaginaries-interpretation} define a
\(K\)-structure on the stated quotient sorts. For every domain assignment
\(\bar a\) and every \(K\)-formula \(\varphi\),
\[
 I(M)\models\varphi(q\bar a)
       \quad\Longleftrightarrow\quad M\models\varphi^I(\bar a).
\]
Every term-value formula selects exactly the representatives of its interpreted
term value. In particular, changing representatives changes no formula truth,
and \(I(M)\models U\) under the stated axiom-translation condition.
No definable selector and no infinite choice of representatives is required.
\end{hbtheorem}

\begin{hbtheorem}{thm:p03-imaginaries-elementarity}{proved=annex:p03-imaginaries-elementarity}{Elementary maps and elementary-chain quotients}
An elementary map \(h:M\to N\) between source models induces an elementary
map \(h^I:I(M)\to I(N)\) given in sort \(s\) by
\(q_s(\bar a)\mapsto q_s(h\bar a)\).
Let \((M_i)_{i<\delta}\) be a nonempty set-indexed elementary chain and
\(M_\infty=\bigcup_{i<\delta}M_i\). The canonical directed limit of
the \(I(M_i)\) is isomorphic to \(I(M_\infty)\), and each canonical
map into that limit is elementary.
\end{hbtheorem}

The word ``canonical'' matters: \([\bar a]_{E_s(M_i)}\) and
\([\bar a]_{E_s(M_j)}\) need not be the same set. The maps above, not
literal inclusion of equivalence classes, identify the successive quotient
structures. This is a genuine difference from the actual-identity
interpretations of \cref{def:p03-vaught-relative-interpretation}.

\begin{hbtheorem}{thm:p03-imaginaries-composition}{proved=annex:p03-imaginaries-composition}{Composition without representative selectors}
Suppose \(J\) interprets an \(H\)-theory in \(T\), and \(I\)
interprets a \(K\)-theory in that \(H\)-theory. Replace each coordinate
of an \(I\)-representative by its \(J\)-representative block. The
flattened domain is the conjunction of those block domains and \(D_s^I\)
translated by \(J\); flattened equality is \((E_s^I)^J\), on that
domain. Translate the relation and function-graph formulas in the same way.
These data define \(I\circ J\), with a canonical isomorphism
\[
                 (I\circ J)(M)\cong I(J(M)).
\]
The formula translation agrees with iterated translation on domain assignments,
up to renaming fresh variables. Dimensions multiply blockwise and add across
different coordinate sorts; quotient equality is not replaced by actual equality.
\end{hbtheorem}

\subsection{Imaginary sorts and canonical parameters}

\begin{hbdefinition}{def:p03-imaginaries-eq-expansion}{The many-sorted imaginary expansion}
For every parameter-free \(L\)-formula \(E\) that \(T\) proves to be
an equivalence relation on a finite product of home sorts, add a distinct sort
\(S_E\), interpreted as that product modulo \(E\), and a total
projection \(q_E\). Retain all home sorts and their original symbols.
Also include each parameter-free definable relation on finite products of these
sorts whose pullback to home representatives is an \(L\)-formula invariant
under the relevant equivalences. This specifies \(L^{\mathrm{eq}}\)
and its structure \(M^{\mathrm{eq}}\); an element \(q_E(\bar a)\)
is an imaginary. Equality on a home sort is one of the equivalences, but an
arbitrary \(S_E\) is not identified with that home sort.
Each formula uses only finitely many sorts and finite representative blocks.
\index{imaginary element}\index{Meq@\(M^{\mathrm{eq}}\)}
\end{hbdefinition}

\begin{hbtheorem}{thm:p03-imaginaries-eq-pullback}{proved=annex:p03-imaginaries-eq-pullback}{Pullback, parameter fibers, and quotients of imaginary sorts}
Every \(L^{\mathrm{eq}}\)-formula pulls back to an \(L\)-formula on
home representatives, with the same truth and elementary preservation.
For a definable \(D\) and equivalence \(E\) on \(D\), extend it to
the whole representative product by
\[
 E^*(\bar x,\bar y)=
 (D(\bar x)\land D(\bar y)\land E(\bar x,\bar y))
 \lor(\neg D(\bar x)\land\neg D(\bar y)\land\bar x=\bar y).
\]
Then \(D/E\) is a definable subset of \(S_{E^*}\). An empty \(D\)
gives an empty definable subset, not an empty home sort.

A parameter-definable equivalence is likewise a fiber of a parameter-free
imaginary sort after the parameters are included as extra representative
coordinates. Finally, a parameter-free definable equivalence on a definable
set of finite tuples from \(M^{\mathrm{eq}}\) has its quotient
definably isomorphic to a definable subset of one of the existing imaginary
sorts. This does not say that such a quotient is a real definable set of \(M\).
\end{hbtheorem}

\begin{hbdefinition}{def:p03-imaginaries-canonical-parameter}{Canonical parameter of a definable family}
For an \(L\)-formula \(\varphi(\bar x,\bar y)\), put
\[
 E_\varphi(\bar y,\bar z)
       =\forall\bar x(\varphi(\bar x,\bar y)
                        \leftrightarrow\varphi(\bar x,\bar z)).
\]
The canonical parameter of \(X=\varphi(M,\bar a)\) is
\(e_X=q_{E_\varphi}(\bar a)\). This parameter belongs to the indicated
imaginary sort, not to an unspecified real tuple. The formula and its
parameter-variable sorts are part of this presentation of the family.
\end{hbdefinition}

\begin{hbtheorem}{thm:p03-imaginaries-canonical-code}{proved=annex:p03-imaginaries-canonical-code}{Canonical parameters and their automorphisms}
The set \(X\) is definable from \(e_X\) in \(M^{\mathrm{eq}}\):
\[
 \bar x\in X\quad\Longleftrightarrow\quad
 \exists\bar y(q_{E_\varphi}(\bar y)=e_X\land\varphi(\bar x,\bar y)).
\]
Every automorphism of \(M\) extends uniquely to the declared
\(M^{\mathrm{eq}}\), and it fixes \(e_X\) exactly when it preserves
\(X\) setwise. These assertions also hold over a named parameter set with
automorphisms fixing that set.
\end{hbtheorem}

\subsection{Real codes, weak elimination, and uniform elimination}

\begin{hbdefinition}{def:p03-imaginaries-real-codes}{Definable and algebraic real codes}
For a parameter set \(A\) in \(M^{\mathrm{eq}}\), an element belongs
to \(\operatorname{dcl}^{\mathrm{eq}}(A)\) if a formula with finitely
many parameters from \(A\) has it as its unique solution in its sort.
It belongs to \(\operatorname{acl}^{\mathrm{eq}}(A)\) if such a
formula has finitely many solutions including it. For tuples these conditions
are coordinatewise, equivalently witnessed by a finite tuple solution set.

A real definable code for \(e\), over \(A\), is a finite tuple
\(\bar c\) of home-sort elements with
\[
 e\in\operatorname{dcl}^{\mathrm{eq}}(A\bar c),\qquad
 \bar c\in\operatorname{dcl}^{\mathrm{eq}}(Ae).
\]
A weak real code replaces the second condition by
\(\bar c\in\operatorname{acl}^{\mathrm{eq}}(Ae)\).
Thus finite ambiguity of an ordered real code is allowed in the weak
condition, but not in interdefinability. Pointwise full elimination means existence
of definable real codes for all imaginaries in the declared models; weak
elimination means existence of weak codes. Neither definition alone asserts
one formula uniform across all models.
\end{hbdefinition}

\begin{hbdefinition}{def:p03-imaginaries-uniform-ei}{Uniform elimination for a theory and a quotient}
For a fixed \(D,E\), uniform elimination over \(T\) means that there
are fixed home-sort output blocks \(\bar z\) and a fixed formula
\(\chi_E(\bar x,\bar z)\) defining, in every \(M\models T\), a
function \(f_E:D(M)\to M^{\bar z}\) with
\[
       E(\bar x,\bar y)\quad\Longleftrightarrow\quad
                         f_E(\bar x)=f_E(\bar y).
\]
Uniform elimination of imaginaries requires this for every declared definable
quotient. A zero-length output block is allowed for a one-class quotient.
The defining formula depends on the equivalence presentation, not on the
choice of model or an individual class. This is the uniformity explicitly
required by Hodges's correction at printed p.~158
\cite[p.~2]{Hodges1993Corrigenda2011}; the lecture presentation also distinguishes
this terminology \cite[pp.~1--2]{Speirs2013Imaginaries}.
\end{hbdefinition}

\begin{hbtheorem}{thm:p03-imaginaries-uniform-codes}{proved=annex:p03-imaginaries-uniform-codes}{Uniform codes, selectors, and the finite-ambiguity boundary}
A uniform \(f_E\) induces a definable bijection from \(D/E\) to its
definable image and gives a definable real code for every class. In particular
it gives a weak code. If a definable selector \(s:D\to D\) satisfies
\(E(\bar x,s(\bar x))\) and is constant on each \(E\)-class, then
\(s\) is a uniform code. Merely choosing one witness separately for each
parameter does not assert that selector's invariance.
Interdefinable codes have the same stabilizer over \(A\); an algebraic
tuple has only a finite orbit under automorphisms fixing \(A\).
\end{hbtheorem}

\begin{hbtheorem}{thm:p03-imaginaries-unordered-pairs}{proved=annex:p03-imaginaries-unordered-pairs}{A weak code that is not a definable real code}
In an infinite pure set \(M\), let \(D=\{(x,y):x\ne y\}\) and let
\(E\) identify the two orderings of the same two-element subset. For
distinct \(a,b\), the imaginary \(e=\{a,b\}\) has weak real code
\((a,b)\) over the empty parameter set, but no definable real code there.
If \(M\) is instead equipped with a linear order, the ordered pair
\((\min(a,b),\max(a,b))\) is a uniform real code for this quotient.
This last conclusion concerns this quotient, not every imaginary of every
ordered theory.
\end{hbtheorem}

\begin{hbtheorem}{thm:p03-imaginaries-infinite-classes}{proved=annex:p03-imaginaries-infinite-classes}{A quotient without even a weak real code}
Let \(M=\mathbb N\times\mathbb N\) in the language of one equivalence
relation \(R\), with \(R((i,j),(k,l))\) iff \(i=k\).
For the class imaginary \(e=q_R((0,0))\), no weak real code over the
empty parameter set exists. A quotient sort still exists: failure of real
coding is not failure of the quotient construction.
\end{hbtheorem}

\subsection{Worked group and field quotients}

\begin{hbtheorem}{thm:p03-imaginaries-quotient-group}{proved=annex:p03-imaginaries-quotient-group}{Interpreting a normal-subgroup quotient}
Let \(G\) be a group and let \(N\) be a definable normal subgroup in
the declared parameter language. On \(D=G\) put
\(E(x,y)\) iff \(x^{-1}y\in N\). The formulas
\[
 G_{\cdot}(x,y,z)=E(xy,z),\qquad
 G_{\mathrm{inv}}(x,z)=E(x^{-1},z),\qquad C_1(z)=E(1,z)
\]
give a quotient interpretation of the group \(G/N\).
Normality is essential to invariance of the multiplication graph: an arbitrary
subgroup gives a coset set, not this quotient group.
\end{hbtheorem}

\begin{hbtheorem}{thm:p03-imaginaries-fraction-field}{proved=annex:p03-imaginaries-fraction-field}{A field interpreted by a two-dimensional quotient}
Let \(R\) be a nonzero commutative integral domain with identity in the
ring language. On \(D=\{(a,b):b\ne0\}\) put
\[
        (a,b)E(c,d)\quad\Longleftrightarrow\quad ad=cb.
\]
The operations on classes are
\[
 [a,b]+[c,d]=[ad+bc,bd],\qquad
 [a,b][c,d]=[ac,bd],\qquad -[a,b]=[-a,b],
\]
with \(0=[0,1]\), \(1=[1,1]\), and inverse \([b,a]\) for a nonzero
\([a,b]\). Their graphs are first-order definable on representatives by
using \(E\) to compare the displayed output to an output variable.
They make \(D/E\) a field and give an injective ring map
\(R\to D/E\), \(a\mapsto[a,1]\).
If \(R\) is already a field, \(f_E(a,b)=z\) defined by \(bz=a\)
is a uniform real code for this quotient. For a general integral domain the
quotient-field construction does not assert a real coding function in \(R\).
\end{hbtheorem}

The public opening of Hodges's original interpretations chapter uses the
two-dimensional construction of complex numbers from real pairs
\cite[Chapter~5, public opening summary]{Hodges1993}. The following local
calculation supplies the field-construction part explicitly, without replacing
the different quotient-field construction above.

\begin{hbtheorem}{thm:p03-imaginaries-complexification}{proved=annex:p03-imaginaries-complexification}{The two-dimensional field construction}
For an ordered field \(F\), use the definable domain \(F^2\), actual
coordinatewise equality, and
\[
 (a,b)+(c,d)=(a+c,b+d),\qquad
 (a,b)(c,d)=(ac-bd,ad+bc),
\]
with zero \((0,0)\), one \((1,0)\), and negation \((-a,-b)\).
These define a field containing the copy \(a\mapsto(a,0)\) of \(F\).
The element \((0,1)\) squares to \((-1,0)\); the inverse of a nonzero
\((a,b)\) is
\[
       \left(\frac{a}{a^2+b^2},\frac{-b}{a^2+b^2}\right),
\]
whose graph is definable by clearing its nonzero denominator. In particular
this gives the usual \(\mathbb C\) from \(\mathbb R\).
The field axioms alone do not establish algebraic closure for an arbitrary
ordered \(F\); no such conclusion is asserted here.
\end{hbtheorem}

\begin{proofboundary}{def:p03-imaginaries-source-boundary}{What a quotient construction does not establish}
Neither adding imaginary sorts nor proving the truth lemma proves elimination
of imaginaries in the original home sorts. The group and field quotients above
do not assert that every group interprets every field, that all theories have
real codes, or that model-dependent codes furnish a uniform theory-level
formula. The original Hodges interpretations presentation includes distinct
relativisation, pseudo-elementary, broad-interpretability and totally
transcendental developments \cite[Chapter~5 contents]{Hodges1993}; these are
not identified with the quotient construction just proved.
\end{proofboundary}
